%% =====================================================================
%%  Chapter 2: Preliminaries
%%  Existing knowledge that this work builds on: the definition of the problem
%%  and the basics of the existing methods. This is NOT your contribution:
%%  summarize textbooks and papers concisely and cite them.
%%  Write your own contribution in the Proposed Method chapter and keep
%%  the two clearly separated.
%%  As an example, this chapter explains the 0-1 knapsack problem and the
%%  greedy algorithm. Replace them with the content of your research.
%% =====================================================================
\chapter{Preliminaries}
\label{chap:preliminaries}

This chapter defines the problem studied in this thesis
and describes the existing method on which the proposed method is based.
\begin{writingnote}
Include only the existing knowledge that the reader needs to understand the proposed method and the experiments;
this chapter is not a textbook.
Explain definitions and existing methods with citations to the original sources (the original papers or standard textbooks).
If a term is defined differently in the literature, state which definition this thesis uses and keep to it.
Follow the notation of the standard textbook in your field and never use one symbol for two meanings;
a list of the main symbols, as in \cref{chap:notation}, helps.
The Academic Phrasebank~\cite{morleyphrasebank} lists phrases commonly used to define terms and to describe existing methods.
\end{writingnote}

\section{Problem Definition}
\label{sec:problem}

%% Describe the problem in words first, then formulate it, then explain the symbols.
%% For multi-line formulas use the align environment; & marks the alignment point.
%% Write the conditions after &&.
%% Refer to consecutive equations with \crefrange{first}{last}, e.g. "Eqs. (2.1) to (2.3)".
As an example, consider the 0-1 knapsack problem~\cite{kellerer2004knapsack}.
Given $n$ items, each with a value and a weight, the problem is to select items
so that the total value is maximized while the total weight does not exceed the capacity.
It is formulated as \crefrange{eq:knapsack-objective}{eq:knapsack-binary}:
\begin{align}
  \text{maximize}   \quad & \sum_{i=1}^{n} v_i x_i \label{eq:knapsack-objective} \\
  \text{subject to} \quad & \sum_{i=1}^{n} w_i x_i \le W, \label{eq:knapsack-capacity} \\
                          & x_i \in \{0, 1\}, && i = 1, \dots, n, \label{eq:knapsack-binary}
\end{align}
where $v_i$ and $w_i$ are the value and the weight of item $i$, $W$ is the capacity,
and $x_i$ is a decision variable that equals 1 if item $i$ is selected and 0 otherwise.
\Cref{eq:knapsack-objective} maximizes the total value of the selected items,
and \cref{eq:knapsack-capacity} ensures that their total weight does not exceed the capacity.

%% When there are only a few symbols, define them in the text right after the
%% equation, as above ("where ...").
%% When there are many symbols (e.g., a large optimization model), define them
%% before the formulation, grouped into sets, parameters and decision variables,
%% with the description environment. Example:
%%   \begin{description}
%%     \item[Sets] $\mathcal{I} = \{1, \dots, n\}$ is the set of items.
%%     \item[Parameters] $v_i$ and $w_i$ are the value and the weight of item $i$,
%%       and $W$ is the capacity.
%%     \item[Decision variables] $x_i$ equals 1 if item $i$ is selected and 0 otherwise.
%%   \end{description}

\section{Existing Method}
\label{sec:greedy}

\begin{writingnote}
Explain the idea and the steps of the existing method so that a reader who does not know it can follow them,
and state its strengths and weaknesses.
Describe the weaknesses concretely so that they lead to the issues and ideas presented in the later chapters.
Write any change that you made to the existing method in the Proposed Method chapter, not here.
Use the present tense to describe existing methods (``The greedy algorithm selects \ldots'').
\end{writingnote}

%% Write pseudocode in the algorithm environment (algorithm2e package). End each line with \;
A representative approximation method for the 0-1 knapsack problem is the greedy algorithm,
which selects items in decreasing order of the ratio $v_i / w_i$~\cite{dantzig1957discrete}.
\Cref{alg:greedy} shows its procedure.
The greedy algorithm is fast, but it does not always find an optimal solution.

\begin{algorithm}[tbp]
  \caption{Greedy algorithm for the 0-1 knapsack problem}
  \label{alg:greedy}
  \KwIn{values $v_1, \dots, v_n$, weights $w_1, \dots, w_n$, capacity $W$}
  \KwOut{set $S$ of selected items}
  Sort the items in decreasing order of $v_i / w_i$\;
  $S \gets \emptyset$, $c \gets 0$ \tcp*{$c$ is the total weight of the selected items}
  \ForEach{item $i$ in the sorted order}{
    \If{$c + w_i \le W$}{
      $S \gets S \cup \{i\}$\;
      $c \gets c + w_i$\;
    }
  }
  \KwRet{$S$}\;
\end{algorithm}
